\section{Introduction}

A finite linear system can preserve every exact state distinction and
nevertheless become progressively less informative under a declared
readout.

That elementary observation becomes structurally nontrivial in the
finite relational mechanics built on
\[
G_{60}\cong \mathrm{AT4val}[60,6].
\]
The established construction already distinguishes quotient-visible
state, lifted state, covering receipt, registered history, additive
winding, and centered scalar magnitude. Projection repeatedly forgets
finite information while more resolved constructions retain distinctions
that are invisible after quotienting or scalar compression.

The present paper asks a narrower question.

Can the exact incidence algebra on
\[
G_{15}\cong L(\mathrm{Petersen})
\]
produce a mathematically controlled concentration of relational
contrast while complete linear distinction remains preserved?

The answer is yes.

Let
\[
M\in\{0,1\}^{15\times30}
\]
be the transport-induced sector-edge incidence matrix and define
\[
Q:=MM^{\mathsf T}.
\]
The established identity
\[
Q=A^3+2A^2+2I
\]
places the sector-overlap algebra inside the adjacency algebra of
\(L(\mathrm{Petersen})\).

The spectrum of \(Q\) is
\[
98^{(1)},\qquad
18^{(5)},\qquad
3^{(4)},\qquad
2^{(5)}.
\]
The eigenvalue \(98\) belongs to the one-dimensional all-ones
common-mode direction. Every centered direction has eigenvalue at
most \(18\).

After the canonical normalization
\[
\widehat Q:=\frac{1}{98}Q,
\]
the common mode is fixed while every centered relational direction
contracts. The sharp centered contraction factor is
\[
\frac{18}{98}=\frac{9}{49}.
\]
Thus
\[
\widehat Q^{\,k}\longrightarrow P_0,
\]
where \(P_0\) is the rank-one orthogonal projection onto the common
mode.

Yet every finite iterate
\[
\widehat Q^{\,k}
\]
is invertible.

This gives an exact finite linear realization of the distinction
between complete state and registered contrast:

\begin{quote}
A strong common aggregate can emerge while relative individuation
becomes arbitrarily weak, without any finite-stage merger of complete
linear states.
\end{quote}

We call this exact algebraic phenomenon
\emph{contrast concentration}.

The broader phrase
\emph{relational saturation} is reserved for a stronger,
projection-relative conjecture. Saturation requires not only a
contrast-concentrating operator, but also an admitted native loading
law, a declared registration, finite independent distinction capacity,
and retained complete information after the declared readout has
ceased to resolve additional independent contrast.

This paper proves contrast concentration.

It formulates relational saturation.

It does not prove that the native higher Thalean mechanics repeatedly
applies \(Q\), nor that physical systems saturate in this manner.
No identification is made here with physical gravity, black holes,
spacetime horizons, Planck-scale information density, or the cosmic
microwave background.

The theorem-side question therefore comes first:

\begin{quote}
Can bounded finite linear aggregation produce strong aggregate and
weak individuation at the same time?
\end{quote}

For the Thalean quadratic, the answer is exact.
